%% kapa: exam=1 round=35 year=2024 period=1 form=A subject=real_estate subject_ko=부동산학원론 seq=15 number=95
%% source: Q-Net 감정평가사 기출문제 (cst00309) artlSeq=5213229 file=raw/1008/5213229/1교시 문제지 원본.pdf
%% conversion: kordoc 4.24.51 -> tools/gen1.py (auto); review=proofread-sonnet; flags=none
\begin{question}{95}
부동산 가격의 제원칙에 관한 내용으로 옳지 않은 것은?
\begin{choices}
  \item 부동산의 가격이 대체ㆍ경쟁관계에 있는 유사한 부동산의 영향을 받아 형성되는 것은 대체의 원칙에 해당된다.
  \item 부동산의 가격이 경쟁을 통해 초과이윤이 없어지고 적합한 가격이 형성되는 것은 경쟁의 원칙에 해당된다.
  \item 부동산의 가격이 부동산을 구성하고 있는 각 요소가 기여하는 정도에 영향을 받아 형성되는 것은 기여의 원칙에 해당된다.
  \item 부동산의 가격이 내부적인 요인에 의하여 긍정적 또는 부정적 영향을 받아 형성되는 것은 적합의 원칙에 해당된다.
  \item 부동산 가격의 제원칙은 최유효이용의 원칙을 상위원칙으로 하나의 체계를 형성하고 있다.
\end{choices}
\end{question}
\answer{4}
